% !PRINTED 37
\weberpage{37}
%<<CONT

\begin{equation}
A=QB+C,
\tag{13}
\end{equation}
\begin{equation}
 C=c_0x+c_1,
\tag{14}
\end{equation}
where
\begin{equation}
 c_0={6a_0a_2-2a_1^2\over 9a_0},
 \qquad
 c_1={9a_0a_3-a_1a_2\over 9a_0}.
\tag{15}
\end{equation}
If $c_0$ is zero, the algorithm is already closed. If $c_1$ is non-zero, then $A$ and $B$ have no common divisor; but if $c_1$ is also zero, then $B$ itself is the greatest common divisor of $A$ and $B$, that is, $A$ is divisible by $B$. The conditions for this are therefore
\begin{equation}
 3a_0a_2-a_1^2=0,
 \qquad
 9a_0a_3-a_1a_2=0.
\tag{16}
\end{equation}
If $c_0$ is not zero, we go one step further and put
\begin{equation}
 B=PC+D,
\tag{17}
\end{equation}
where $D$ becomes constant and has the expression
\begin{equation}
 D={a_2c_0^2-2a_1c_0c_1+3a_0c_1^2\over c_0^2}.
\tag{18}
\end{equation}
If this expression is non-zero, then $A$ and $B$ are \gesperrt{prime to one another}; if it is zero, then $A$ and $B$ have the greatest common divisor $C$. If for $c_0,c_1$ one substitutes the values (15), omits the denominator, still removes the non-vanishing factor $9a_0$, and changes the sign, this condition, after a simple calculation, assumes the form
\begin{equation}
 a_1^2a_2^2+18a_0a_1a_2a_3-4a_0a_2^3-4a_1^3a_3-27a_0^2a_3^2=0.
\tag{19}
\end{equation}
It is, as one readily sees either by calculation or from (18), also fulfilled when the conditions (16) hold, and is therefore the necessary and sufficient condition that $f(x)$ and $f'(x)$ have a common divisor. The left hand side of (19), which is a whole rational and homogeneous function of the coefficients of $f(x)$, is called the \gesperrt{discriminant of the function $f(x)$}.

We derive one more theorem from the algorithm (3), which is often used.

From the first of these equations it follows that
\begin{equation}
 A''=A-Q'A',
\tag{20}
\end{equation}
and if this value of $A''$ is substituted in the following equation, one obtains
\[
 A'''=(1+Q'Q'')A'-Q''A,
\]
%CONT>>
